\section{Baseline walkers and the exact kernel of a sweep}
\label{sec:c-baseline-forest}

For this method it is enough to bound the relative entropy of a large
marginal by a power of $\log n$; it need not tend to zero.  A walker
moving along the paths of one realization of the shuffle gives that
weaker bound by a collision calculation.  We then obtain a separate
estimate for representations with a long first row directly from the
unique paths through one sweep.  The two estimates enter different
ranges of the abstract hybrid transfer. Unlike the vanishing entropy
estimate of Section~\ref{lifts:cycle-information}, this is an alternative
derivation that combines weak entropy with additional sparse Fourier
information.

Throughout this section $V=\mathbb F_2^d$, $n=2^d$, and layers use the
coordinate directions $1,\ldots,d$ cyclically, with independent fair
switches.  One sweep starts at direction $1$ and consists of $d$
physical shuffles.  Write $K$ for its permutation law.  All logarithms
in entropy expressions are natural logarithms.

\subsection{A walker on a baseline realization}

\begin{proposition}[A weak entropy bound for large marginals]
\label{prop:c-baseline-entropy}
There is an absolute positive integer $C_0$ such that the following
holds for all sufficiently large $d\ge4$.  Let $\mu$ be the permutation
law after $C_0d$ layers.  If $I$
is a uniform ordered list of $15n/16$ distinct input labels, independent
of the shuffle, and $u_I$ is its uniform injection law, then
\[
 \mathbb E_I D(\mu_I\Vert u_I)=O((1+d)^4).
\]
\end{proposition}

We first prove the collision estimate that supplies the proposition.
Fix a distinguished input label $j$ and a uniform set $A$ of $m$ active
labels containing $j$, where $m\ge\epsilon n$ and $\epsilon=1/16$.
The labels outside $A$ will be revealed.  Independently sample a
\emph{baseline} realization of every switch.  A second walker starts
at $j$ and travels along the baseline paths.  When its current baseline
path is paired with another active path, toss a fresh fair coin to
decide whether to toggle to that other path.  If the partner is
inactive, stay on the current path.  Staying means following that
baseline path through the current switch.

Given all revealed inactive paths, let $q$ be the conditional endpoint
law of $j$. Lemma~\ref{lem:marginal-averaging} leaves every unvisited
switch independent and fair. Fix the entire baseline realization and
average only the extra toggle coins. The walker's endpoint still has
law $q$: a pair of active paths
offers the two available outputs with equal probabilities, and an
active--inactive pair has only one available output. These are exactly
the conditional-flow rules, independently of the baseline choices on
active--active switches. If $X_j$ is the baseline endpoint of $j$,
the conditional collision probability given that baseline is therefore
$q(X_j)$. Averaging first over baselines with the same inactive paths
gives $\sum_xq(x)^2$, and hence
\begin{equation}
 \mathbb E\sum_xq(x)^2
 =\mathbb P\{\hbox{walker and baseline }j
                    \hbox{ have the same endpoint}\}.
 \label{eq:c-baseline-collision-identity}
\end{equation}
The expectation includes the random active set and baseline.

\begin{lemma}[Coincidence estimate]
\label{lem:c-baseline-coincidence}
For a sufficiently large absolute integer $C_0$, at $t=C_0d$,
\begin{equation}
 \mathbb P\{\hbox{walker and baseline }j\hbox{ coincide at }t\}
 \le \frac1m+O\bigl((1+t)^3/n^2\bigr),
 \label{eq:c-baseline-coincidence}
\end{equation}
uniformly over $m\ge\epsilon n$ and $j$.
\end{lemma}

\begin{proof}
We first compare the xor of the two positions with an ideal chain
that contracts over blocks of four sweeps. We then control the actual
transition error by estimating coincidence and the joint event of
coincidence with a repeated partner. Two inverse-$n$ costs make this
error small enough to telescope through the remaining kernels.

Write $D_s$ for the binary xor of the walker position and the position
of the baseline path of $j$ immediately before layer $s$.  We generate
the baseline chronologically, and reveal whether a label is active
only when the walker queries it.  Apart from $j$, activity statuses
are a uniform sample of $m-1$ active labels from $n-1$ possibilities.
A partner's identity is known before the current layer's fresh coins
are sampled, so such adaptive queries preserve sampling without
replacement among unqueried labels.  At most $1+t$ labels are queried
through time $t$.

Consider layer $s$ in direction $i=i_s$.  If
$D_s\notin\{0,e_i\}$, the two distinguished paths use distinct
switches.  Their new xor bit in coordinate $i$ is fair, independently
of the history, and the other bits are unchanged.  If $D_s=e_i$, the
walker is paired with $j$ and a fair toggle makes the new xor either
$0$ or $e_i$.  If $D_s=0$, the new xor becomes $e_i$ precisely when
the partner of $j$ is active and the walker toggles to it.

Define ideal kernels $Q_s$ by the same rules, replacing the last
activity probability by the constant
\[
 p=\frac{m-1}{n-1}.
\]
They have a common stationary law
\[
 \pi(0)=\frac1m,\qquad \pi(x)=\frac p m\quad(x\ne0).
\]
Indeed $1+(n-1)p=m$.  The edge $\{0,e_i\}$ balances because
$\pi(0)p/2=\pi(e_i)/2$, and every other coordinate pair has equal
masses and a fair refresh.

We need a uniform contraction estimate for products of these ideal
kernels.  Call a layer exceptional when its input is $0$ or $e_i$.
Once a run of consecutive exceptional layers ends, its output is a
singleton.  That singleton bit cannot disappear until its coordinate
is next updated, so there are at least $d-1$ nonexceptional layers
before a new exceptional run.  A new run following these layers
requires $d-1$ specified zero outputs of fair refreshed bits and has
probability at most $2^{-(d-1)}$.  A run continuing for at least $d$
layers must keep outputting zero until its last layer; from any
specified start its probability is at most
$(1-p/2)^{d-2}$.  This also covers the first singleton-to-zero move,
whose probability $1/2$ is no larger than $1-p/2$.

In a block of $4d$ layers, a union bound therefore gives, uniformly in
the starting state,
\[
 \mathbb P\{\hbox{an exceptional layer occurs in the last }d
                   \hbox{ layers}\}
 \le C d\bigl(2^{-(d-1)}+(1-p/2)^{d-2}\bigr).
\]
A run that starts early and reaches the last $d$ layers is long;
a later start has the preceding $d-1$ fair refreshes just described.
Couple two ideal chains by common random bits, using the same refreshed
bit whenever both are nonexceptional and keeping them together once
they agree.  If neither has an exception in the last $d$ layers, all
coordinates have been refreshed identically and the endpoints agree.
Since $p\ge\epsilon/2$ for large $n$, there is an absolute $\beta>0$
for which every $4d$-layer ideal product has total-variation contraction
coefficient at most $n^{-\beta}$.

We now compare the actual xor law $\gamma_s$ with this ideal chain.
Only the case $D_s=0$ can differ.  At an unqueried partner, sampling
without replacement changes its conditional activity probability from
$p$ by $O((1+t)/n)$, uniformly through time $t=O(d)$.  At a previously
queried partner use the bound one.  Thus
\begin{equation}
 \|\gamma_{s+1}-\gamma_sQ_s\|_{\rm TV}
 \le\mathbb P\{D_s=0,\ \hbox{partner previously queried}\}
       +O((1+t)/n)\mathbb P\{D_s=0\}.
 \label{eq:c-baseline-one-step-error}
\end{equation}
The remaining task is to get a second factor $1/n$ in the first term.

First dispose of trajectories that have never split from $j$. Define
an auxiliary test along the baseline path of $j$, regardless of whether
the actual walker has already left it. Sample an independent potential
toggle coin at every layer, using it for the actual walk whenever a
toggle is allowed. In each successive block of
$2d$ layers, monitor only the last $d$ layers. At a monitored layer,
query the partner's activity and use the walker's toggle coin to test
for an active toggle. Reveal no unmonitored activity statuses for this
test. If the actual walker never splits, every monitored test fails.
We bound this larger event using only the monitored-test history.
The partners of $j$ at these monitored layers are distinct: once a
baseline path is paired with $j$, it leaves with an opposite output
bit, and that bit is not updated again during the monitored interval.
Condition on the baseline history and monitored queries at the
beginning of the block. For any
previously queried label and any monitored layer, becoming the partner
of $j$ requires matching the preceding $d-1$ output bits while using
distinct switches.  Otherwise a previous shared switch leaves an
opposite bit that persists.  Sequential exposure of the fresh distinct
switches bounds this probability by $2^{-(d-1)}$.
A union bound over the $d$ layers and $O(1+t)$ known labels bounds the
probability of an old partner by $O(d(1+t)/n)$.

On a sequence of fresh partner queries the remaining active density
is at least $\epsilon/2$ for large $n$, since $t=O(d)$.  At each
monitored layer an active toggle therefore has conditional probability
at least $\epsilon/4$.  The probability of avoiding all such toggles
while the partners stay fresh is at most $(1-\epsilon/4)^d$.
Repeating this conditional estimate at successive block starts, with
no conditioning on unmonitored toggle outcomes, gives
\[
 \mathbb P\{\hbox{no split through }C_1d\}
 \le\bigl(O(d(1+t)/n)+(1-\epsilon/4)^d\bigr)^{\lfloor C_1/2\rfloor}
 \le n^{-3}
\]
for a sufficiently large absolute integer $C_1$.

Any other coincidence at time $s\ge C_1d$ has a last collapse layer
$u<s$: its input xor is $e_{i_u}$, it collapses to zero, and it stays
zero through time $s$.  Necessarily $u\ge d$, since the singleton
created at the initial split persists until its next coordinate
update.  In the $d-1$ layers $u-d+1,\ldots,u-1$, the walker and $j$
use distinct switches and output equal updated bits.  If they were
coincident in that interval, a subsequent split would leave a bit
that could not disappear before $u$; an earlier partner layer likewise
cannot produce a different admissible history.  Consequently the
contribution from a fixed $u$ is at most $2^{-(d-1)}$, and
\begin{equation}
 \mathbb P\{D_s=0\}\le n^{-3}+2t/n.
 \label{eq:c-baseline-rough-coincidence}
\end{equation}

For the joint event in \eqref{eq:c-baseline-one-step-error}, fix a
last collapse layer $u$ and let $E_u$ be the event that the
walker collapses at $u$, stays on $j$ through time $s$, and the partner
at $s$ has been queried previously. The first $1/n$ cost comes from
the equal-bit window ending at $u$; the second comes from the final
partner matching $j$ in the window ending at $s$. These windows may
overlap, so their costs cannot be multiplied by an independence claim.

Start from the true history at $u-d+1$ and define a sequential law
$\widetilde{\mathbb P}_u$ that is allowed to abort. At each of the next
$d-1$ layers, query the
needed status and choose the walker toggle as usual.  Abort if the two
distinguished switches coincide.  Otherwise sample their fresh coins
conditional on the resulting output bits being equal; this conditioning
has probability exactly $1/2$.  Sample all other switches ordinarily.
At layer $u$, abort unless the two paths are partners, then force the
walker to collapse onto $j$.  Thereafter query the needed statuses and
force the walker to stay on $j$ through time $s$.  No conditioning on
a future successful completion is used. Every true history contributing
to $E_u$ has likelihood at most $2^{-(d-1)}$ times its likelihood under
$\widetilde{\mathbb P}_u$: the equalities each cost $1/2$, and the forced
decisions discard only factors at most one. It remains to bound the
probability of a previously queried final partner under this law.

Let $W$ consist of the $d-1$ layers immediately preceding $s$.
Throughout $W$ the two distinguished outputs have matching updated
bits, either by the equality construction or by coincidence after $u$.
A baseline label $b\ne j$ that is the partner of $j$ at $s$ cannot
have touched either distinguished switch during $W$.  At a switch
of $j$ it would leave with the opposite bit.  At a switch of the
walker it either leaves with the opposite bit or follows the walker;
in the latter case it must later leave before time $s$, again creating
an opposite bit.  Every coordinate used in $W$ is different from
$i_s$ and is used only once, so any such bit would persist and preclude
partnership at $s$.  It follows also that a previously queried final
partner was queried before $W$: all status queries inside $W$ concern
labels touching the walker's switch there.

Condition on the simulated history before $W$.  There are only
$O(1+t)$ earlier queried candidates.  For each candidate, matching the
$d-1$ required bits without touching either distinguished switch has
probability at most $2^{-(d-1)}$.  At each step its separate switch is
fresh and independent of the distinguished bit, even when the two
distinguished coins are being sampled conditional on equality.  Abort
conditions can only reduce this probability. Thus the probability of a
previously queried final partner under $\widetilde{\mathbb P}_u$ is
$O((1+t)/n)$. The likelihood comparison now gives
\[
 \mathbb P(E_u)\le 2^{-(d-1)}O((1+t)/n)
                    =O((1+t)/n^2).
\]

Sum over $u<s$ and include the no-split contribution.  Together with
\eqref{eq:c-baseline-rough-coincidence}, this proves, for $s\ge C_1d$,
\[
 \|\gamma_{s+1}-\gamma_sQ_s\|_{\rm TV}
 =O((1+t)^2/n^2).
\]
Now follow $C_1d$ by a fixed number $b>3/\beta$ of $4d$-layer blocks.
The ideal product sends any law to within $n^{-3}$ of $\pi$.
Telescoping the actual errors through the remaining Markov kernels,
which contract total variation, costs $O((1+t)^3/n^2)$.
With $C_0=C_1+4b$ we have
\[
 \|\gamma_t-\pi\|_{\rm TV}
 \le n^{-3}+O((1+t)^3/n^2),\qquad t=C_0d.
\]
Since $\pi(0)=1/m$, this proves
\eqref{eq:c-baseline-coincidence}.
\end{proof}

\begin{proof}[Proof of Proposition~\ref{prop:c-baseline-entropy}]
Use the entropy chain rule for the uniformly ordered input list.
At each rank, reveal in addition the entire paths of the earlier
labels; this can only increase the conditional relative entropy from
the uniform available endpoint.  The remaining active set has size
$m\ge n/16$, and its next label is uniform in that set.  By label
symmetry of the uniform order, one may fix that next label as $j$ and
take the other active labels uniformly among the remaining labels,
exactly as in Lemma~\ref{lem:c-baseline-coincidence}.
For the resulting conditional endpoint law $q$,
\[
 D(q\Vert U_m)\le m\sum_xq(x)^2-1.
\]
Equations~\eqref{eq:c-baseline-collision-identity} and
\eqref{eq:c-baseline-coincidence} bound its mean by
$O(m(1+t)^3/n^2)$.  Summing over at most $n$ ranks gives
$O((1+t)^3)$, and $t=C_0d$ proves the stated result.
\end{proof}

\subsection{The exact kernel on a short ordered tuple}

The weak entropy estimate leaves representations with a long first
row for a different argument.  Their action is visible in short tuples,
where a sweep has an explicit endpoint kernel.
For a partition $\lambda\vdash n$, let $D_\lambda$ be its irreducible
degree and let $\widehat K(\lambda)$ be the average of its unitary
representation matrices under $K$.

The next lemma isolates the path calculation common to this operator
estimate and the Hilbert--Schmidt estimate in
Section~\ref{c:sec:replica}. The probabilistic entropy arguments in the
two sections are different; only this one-sweep kernel is shared.

\begin{lemma}[The prescribed-path kernel and isolation cancellation]
\label{lem:c-sweep-kernel}
Let $1\le k\le n$ and let $x=(x_a)_{a=1}^k$ and
$y=(y_a)_{a=1}^k$ be tuples distinct on each side. The prescribed path
from $x_a$ to $y_a$ replaces the input bits by the output bits in
sweep order. For $a<b$, let $J$ be the last
coordinate in which $x_a,x_b$ differ and $I$ the first coordinate in
which $y_a,y_b$ differ.  Define
\[
 u_{ab}=\begin{cases}
 0,&I<J,\\
 1,&I=J,\\
 -1,&I>J.
 \end{cases}
\]
For $T\subseteq[k]$, put
\[
 F_T(x,y)=\prod_{a<b\,;\ a,b\in T}(1+u_{ab}),\qquad
 F_*(x,y)=\sum_{T\subseteq[k]}(-1)^{k-|T|}F_T(x,y).
\]
Then $F_T=n^{|T|}\Pr_K\{g x_a=y_a\text{ for every }a\in T\}$.
In particular, the tuple transition kernel is
\begin{equation}
 K(x,y)=n^{-k}F_{[k]}(x,y).
 \label{eq:c-true-sweep-pair-kernel}
\end{equation}
Join $a,b$ when their prescribed paths use a common switch; equivalently,
when $u_{ab}\ne0$. If this graph has an isolated vertex, then $F_*=0$.
For all such tuples,
\begin{equation}
 |F_*(x,y)|\le 2^k k^{k/2}.
 \label{eq:c-true-sweep-prefix-bound}
\end{equation}
\end{lemma}

\begin{proof}
The prescribed path of card $a$ is unique: successively replace its
input bits by its output bits in sweep order. The three cases in the
definition of $u_{ab}$ describe its interaction with the path of $b$.
The first case has no shared switch, the second has exactly one shared
switch with consistent opposite outputs, and the third forces a vertex
collision. For $r=|T|$ prescribed paths, the bits cost $2^{-dr}$
before identifying shared switches; each consistent shared switch saves
a factor two. A collision makes both the probability and the product
zero. This proves the formula for $F_T$. If a vertex is isolated,
including or omitting it does not change $F_T$, so its two terms in the
alternating sum cancel.

To bound the alternating sum, consider a compatible subset $T$ of
$r\le k$ paths. Its product is $F_T=2^{s_T}$, where $s_T$ counts
their shared switches. The empty subset contributes one; for $r\ge1$,
arrange the output words in a
binary prefix tree.  At a node with child counts $a,b$, every shared
switch sends one path into each child, so at most $\min(a,b)$ switches
are shared there.  The inequality
\[
 \min(a,b)\le\frac{(a+b)\log_2(a+b)-a\log_2a-b\log_2b}{2}
\]
uses $0\log_2 0=0$ and follows from concavity of binary entropy,
which is at least twice its smaller argument.  Sum over all nodes;
the right side telescopes to $r\log_2r/2$, bounding $s_T$. The product is at most
$r^{r/2}\le k^{k/2}$, and summing its at most $2^k$ appearances proves
\eqref{eq:c-true-sweep-prefix-bound}.
\end{proof}

\begin{proposition}[True-sweep sparse estimate and row cutoff]
\label{prop:c-true-sweep-sparse}
For all sufficiently large $d$,
\begin{equation}
 \|\widehat K(\lambda)\|_{\rm op}\le n^{-k/10}
 \quad\hbox{if}\quad 1\le k=n-\lambda_1\le n^{1/20}.
 \label{eq:c-true-sweep-sparse}
\end{equation}
If $\lambda$ has more than $n/2$ rows, then
$\widehat K(\lambda)=0$.
\end{proposition}

\begin{proof}
We use the classical branching and Young rules in the conventions of
\cite{sagan,etingof}. The permutation representation on ordered distinct
$k$-tuples is $\operatorname{Ind}_{S_{n-k}}^{S_n}\mathbf1$.
By branching it contains $\lambda$ precisely when $\lambda_1\ge n-k$.
When $k=n-\lambda_1$, its $\lambda$-isotypic space is orthogonal to
functions omitting any specified coordinate: these carry the
$(k-1)$-tuple representation, which contains no copy of $\lambda$.
Thus each function in this space, extended by zero to tuples with
repetitions, has sum zero on varying any one coordinate.

Take two such test functions and use the kernel of
Lemma~\ref{lem:c-sweep-kernel}. Extend each $u_{ab}$ arbitrarily on
pair repetitions, preserving its dependence on just the two indicated
input and output coordinates. Every proper-subset term $F_T$ vanishes
in the bilinear form after integrating an omitted coordinate, so we
may replace $F_{[k]}$ by $F_*$.
Under independent uniform input and output coordinates, each
zero-extended function has squared norm $a_k=(n)_k/n^k$ times its
squared norm under the uniform injection law. By
\eqref{eq:c-true-sweep-pair-kernel}, the bilinear form has the same
factor $a_k$ relative to the Markov bilinear form. Hence the resulting
operator bound transfers without loss to the injection normalization.
Acting on functions instead of masses may replace the Fourier matrix
by its adjoint, which has the same operator norm.

If there are no isolated vertices, choose a spanning forest in each
nontrivial component of the nonzero-edge graph.  Its number of edges
satisfies $k/2\le e\le k-1$.  For an upper bound there are at most
$k^{2e}$ candidate forests with $e$ edges and at most $d^e$ choices of
the edge coordinates $J$.  A nonzero edge requires equality of the
input suffix after $J$ and equality of the output prefix before $J$.
For a fixed forest and choices of $J$, write its requirement indicator
as $A(x)B(y)$.  For each bit, the edges imposing equality form a
subforest, so there is one independent binary constraint per edge;
choose the bits successively from roots to leaves.  It follows that
\[
 \mathbb E_x A(x)=2^{-\sum_{\rm edges}(d-J)},\qquad
 \mathbb E_y B(y)=2^{-\sum_{\rm edges}(J-1)}.
\]
The kernel $A(x)B(y)$ has operator norm equal to the geometric mean
of these probabilities, namely $(2/n)^{e/2}$.  One may also obtain the
same value by the row and column bounds in Schur's test.
Taking absolute values of the test functions, using
\eqref{eq:c-true-sweep-prefix-bound}, and summing this domination gives
for $k\ge2$ the bound
\[
 2^k k^{k/2}
   \sum_{\lceil k/2\rceil\le e\le k}
       \bigl(k^2d\sqrt{2/n}\bigr)^e.
\]
It is at most $n^{-k/10}$ uniformly for $2\le k\le n^{1/20}$ and
large $n$.  For example,
$2^k k^{k/2}\le n^{k/20}$ and
$k^2d\sqrt{2/n}\le n^{-3/8}$ eventually, so the displayed sum is
at most $2n^{-11k/80}\le n^{-k/10}$.  When $k=1$, $F_*=1-1=0$,
which is the exact one-card uniformity of a sweep.  This proves
\eqref{eq:c-true-sweep-sparse}.

Finally a single switch layer is averaging over the subgroup
$(S_2)^{n/2}$ supported on its disjoint matching pairs.  Young's rule
says its invariant vectors can occur only in shapes dominating
$(2^{n/2})$.  Such a shape has at most $n/2$ rows, since the sum of its
first $n/2$ rows must be at least $n$.  Thus every layer projection,
and hence the sweep product, vanishes on all taller shapes.
\end{proof}

\subsection{Applying the entropy--sparse hybrid transfer}

\begin{corollary}[Baseline-walker application]
\label{cor:c-baseline-forest-mixing}
For the absolute integer $C_0$ in
Proposition~\ref{prop:c-baseline-entropy}, the full permutation law
after $(12C_0+40)d$ physical shuffles tends to uniform in total
variation, uniformly over initial decks.
\end{corollary}

\begin{proof}
Choose a uniform partition of the input labels into sixteen equal
blocks.  Each block complement has the law of the retained domain in
Proposition~\ref{prop:c-baseline-entropy}; ordering that complement
does not affect its relative entropy.  Averaging the sum over all
sixteen complements therefore selects a deterministic partition with
total marginal relative entropy $O((1+d)^4)$.

The entropy and sparse-degree hybrid theorem
\cite[Theorem~\ref*{l-l:forest-hybrid}]{partialFourier}
applies to this fixed partition and to the
law $\mu$ of $C_0$ sweeps.  Its entropy truncation deletes permutations
whose density on any complementary marginal exceeds
$\exp(n^{1/50})$, then normalizes to a law $\nu$, as in
\cite[Lemma~\ref*{l-l:forest-clipping}]{partialFourier}.  To check its mass
hypothesis directly, for a density $f$ relative to a probability $u$,
\[
 \int_{\{f>e^a\}}f\,du
 \le\frac{D(fu\Vert u)+1/e}{a},
\]
because the negative part of $f\log f$ is at most $1/e$ pointwise.
Taking $a=n^{1/50}$ and summing over the sixteen marginals gives
$\|\nu-\mu\|_{\rm TV}=O((1+d)^4/n^{1/50})=o(1)$.
After normalization each complementary coset density is at most
$2\exp(n^{1/50})$ for large $n$.

The large-degree estimate
\cite[Proposition~\ref*{l-l:forest-bulk}]{partialFourier} gives
$\|\widehat\nu(\lambda)\|_{\rm op}\le D_\lambda^{-1/3}$ whenever
$k=n-\lambda_1>n^{1/20}$ and $\lambda$ has at most $n/2$ rows.
The hybrid theorem
\cite[Theorem~\ref*{l-l:forest-hybrid}]{partialFourier} also assumes a
probability law $K$ with operator decay $n^{-k/10}$ for
$1\le k\le n^{1/20}$ and zero transform on taller shapes.
Proposition~\ref{prop:c-true-sweep-sparse} supplies these estimates for
one true sweep, separately from the entropy construction.

The theorem controls the exact convolution
$\sigma=\nu^{*12}*K^{*40}$, whose Fourier matrix in our convention is
$\widehat\sigma(\lambda)=\widehat\nu(\lambda)^{12}
\widehat K(\lambda)^{40}$. All convolution factors are independent.
Its operator norm is at most $D_\lambda^{-4}$ in the large-deficit
range, at most $n^{-4k}$ in the small-deficit range, and zero above
$n/2$ rows.  The Plancherel sum uses
$D_\lambda^2$ times the square of these operator bounds, with ordinary
unnormalized Hilbert--Schmidt norms.  The degree and partition estimates
in \cite[Theorem~\ref*{l-l:forest-hybrid}]{partialFourier}
sum the large range to $o(1)$; in the small range
$D_\lambda\le n^k$ from the tuple module, so the contribution at each
$k$ is at most the number of partitions of $k$ times $n^{-6k}$, whose
sum also tends to zero.  The sign shape has more than $n/2$ rows and
is annihilated by $K$, so no parity bias survives.  Cauchy--Schwarz
therefore gives total variation tending to zero.

Replacing the twelve factors $\nu$ by $\mu$ costs at most
$12\|\nu-\mu\|_{\rm TV}=o(1)$.  The resulting law consists of
$12C_0+40$ successive true sweeps, which is $(12C_0+40)d$ physical
shuffles.  Translation by an initial deck preserves the same distance.
\end{proof}
